\documentclass[11pt]{amsart}
\usepackage{amsmath,amssymb,amsthm}
\usepackage{booktabs}
\usepackage{xcolor}
\usepackage[no-math]{fontspec}
\setmonofont{DejaVuSansMono.ttf}[Scale=0.82]
\usepackage[margin=1.1in]{geometry}
\usepackage[colorlinks=true,linkcolor=blue!50!black,citecolor=blue!50!black,urlcolor=blue!50!black]{hyperref}
\usepackage{fancyvrb}

\newtheorem{theorem}{Theorem}[section]
\newtheorem{lemma}[theorem]{Lemma}
\newtheorem{corollary}[theorem]{Corollary}
\newtheorem{proposition}[theorem]{Proposition}
\newtheorem*{dpcor}{Corollary~2 of \cite{DP}}
\newtheorem*{dpconj}{Conjecture~1 of \cite{DP}}
\theoremstyle{remark}
\newtheorem{remark}[theorem]{Remark}

\newcommand{\cE}{\mathcal{E}}
\DeclareMathOperator{\tr}{tr}
\DeclareMathOperator{\Cay}{Cay}

\title[Locally equienergetic graphs]{Connected locally equienergetic graphs:\\ counterexamples to a conjecture of Dede and Popat}

\author{GPT-6-Astra}
\author{Claude Fable 5.1}
\author{Claude Opus 5.5}
\thanks{The authors are AI models: GPT-6-Astra (OpenAI) and Claude Fable~5.1 and Claude Opus~5.5 (Anthropic). Source files, scripts and Lean code: \url{https://github.com/Kevin-OS7/math}.}

\date{October 29, 2026}

\begin{document}

\begin{abstract}
The energy of a graph is the sum of the absolute values of the eigenvalues of its adjacency matrix.
The local energy of a graph is the sum, over all vertices $v$, of the decrease of the energy when $v$ is deleted.
Dede and Popat called two graphs of the same order locally equienergetic if their local energies are equal.
On the basis of a computer search they stated that, up to order $10$, the only pairs of connected locally equienergetic graphs are $(K_n,C_n)$ with $n$ odd, and they conjectured that there are no other pairs at all (MATCH Commun.\ Math.\ Comput.\ Chem.\ 93 (2025), Corollary~2 and Conjecture~1).
We show that both statements are false.
Among connected graphs of order at most $10$ there are exactly $58$ further pairs, all of order $10$.
In each of them the two graphs have the same eigenvalues and angles, and they are related by a Godsil--McKay switching.
At order $16$, the $4\times4$ rook's graph, the Shrikhande graph and the complement of the Clebsch graph all have local energy $64$, although the energy of the last graph is $40$ and that of the other two is $36$.
More generally, cospectral walk-regular graphs are locally equienergetic, and a strongly regular graph of order $n$ whose eigenvalues $k>r\ge0>s$ satisfy a sign condition has local energy $2n|s|$.
This gives infinite families of counterexamples, including pairs of graphs with different spectra.
The counterexamples of orders $10$ and $16$ are verified in Lean~4.
\end{abstract}

\maketitle

\section{Introduction}

Let $G$ be a finite simple graph with $n$ vertices and adjacency matrix $A(G)$.
The energy of $G$ is $\cE(G)=\sum_{i=1}^n|\lambda_i|$, where $\lambda_1,\dots,\lambda_n$ are the eigenvalues of $A(G)$; see \cite{LSG} for background.
For a vertex $v$, let $G-v$ denote the graph obtained from $G$ by deleting $v$.
Espinal and Rada \cite{ER} studied the difference $\cE(G)-\cE(G-v)$, which is called the local energy of $G$ at $v$.
Following \cite{ER,DP}, the local energy of $G$ is
\[
e(G)=\sum_{v\in V(G)}\bigl(\cE(G)-\cE(G-v)\bigr).
\]
Dede and Popat \cite{DP} call two graphs of the same order \emph{locally equienergetic} if their local energies are equal.
As recalled in \cite{DP}, formulas of Espinal and Rada give $e(K_n)=2n$ for $n\ge2$ and $e(C_n)=2n$ for odd $n$.
Hence $K_n$ and $C_n$ are locally equienergetic for odd $n$.
Dede and Popat report an exhaustive computer search over the graphs of order at most $10$, list the classes of locally equienergetic graphs of orders $4$ to $8$, and state the following.

\begin{dpcor}
There exists no simple, undirected, and connected locally equienergetic graph of order $n\le10$ except the pair $(K_n,C_n)$ for odd $n$.
\end{dpcor}

\begin{dpconj}
There exists no simple, undirected, and connected locally equienergetic graph other than the pair of locally equienergetic graphs $(K_n,C_n)$ for odd $n$.
\end{dpconj}

Every graph is locally equienergetic with any graph isomorphic to it.
We therefore read both statements as statements about pairs of non-isomorphic graphs: if $G_1$ and $G_2$ are non-isomorphic connected graphs of the same order $n$ with $e(G_1)=e(G_2)$, then $n$ is odd and $\{G_1,G_2\}=\{K_n,C_n\}$ up to isomorphism.
In Corollary~2 the order is restricted to $n\le10$.
The definition of local equienergy in \cite{DP} does not exclude cospectral pairs.
We show that both statements are false, and that the conclusion of Conjecture~1 also fails for pairs of graphs with different energies.

\begin{theorem}\label{thm:10}
Let $n\le10$.
The pairs of non-isomorphic connected locally equienergetic graphs of order $n$ are $(K_n,C_n)$ for $n\in\{5,7,9\}$ and $58$ pairs of order $10$.
In each of the $58$ pairs the two graphs have the same eigenvalues and angles, the same degree sequence and cospectral complements.
In each pair, one graph is isomorphic to the graph obtained from the other by Godsil--McKay switching with respect to a set of four vertices.
\end{theorem}

Theorem~\ref{thm:10} is the result of a computer search, described in Section~\ref{sec:search}.
It shows that Corollary~2 of \cite{DP} is false and that the smallest counterexample to Conjecture~1 has order $10$.
The pair of order $10$ with the fewest edges is given explicitly in Section~\ref{sec:ten}, and the fact that it is a counterexample is verified in Lean.

\begin{theorem}\label{thm:16}
Let $R$ be the $4\times4$ rook's graph $K_4\square K_4$, let $S$ be the Shrikhande graph, and let $Q$ be the complement of the Clebsch graph.
These three graphs are connected, have $16$ vertices and are pairwise non-isomorphic.
They satisfy $e(R)=e(S)=e(Q)=64$, $\cE(R)=\cE(S)=36$ and $\cE(Q)=40$.
\end{theorem}

The pair $(R,S)$ is a pair of cospectral graphs.
The pair $(R,Q)$ is not even equienergetic.
So Conjecture~1 also fails if one requires, as is usual for equienergetic graphs, that the two graphs are not cospectral.

Both examples are instances of two general mechanisms.
Cospectral walk-regular graphs have the same eigenvalues and angles and are therefore locally equienergetic (Corollary~\ref{cor:walk}).
A connected strongly regular graph of order $n$ with eigenvalues $k>r\ge0>s$ and multiplicities $1,f,g$ has local energy $2n|s|$ whenever $rs+k(fs+gr)\ge0$ (Proposition~\ref{prop:srg}).
For example, the triangular graph $T(m)$ and the lattice graph $L_2(\ell)$ are locally equienergetic whenever $m(m-1)/2=\ell^2$ (Corollary~\ref{cor:fam}).
This equation has infinitely many solutions, and the two graphs have different degrees.

The following are verified in Lean~4 with Mathlib (Section~\ref{sec:lean}): the energies and local energies in Theorem~\ref{thm:16}, the connectivity of $R$, $S$ and $Q$, the facts $R\not\cong S$ and $R\not\cong Q$, the refutation of Conjecture~1 by the pairs $(R,S)$ and $(R,Q)$, and the refutation of Corollary~2 and Conjecture~1 by the explicit pair of order $10$.

\section{Vertex-deleted subgraphs and angles}\label{sec:angles}

Let $\phi(G,x)=\det(xI-A(G))$ be the characteristic polynomial.
For a polynomial $p$ with only real roots, let $\cE(p)$ be the sum of the absolute values of its roots, counted with multiplicity.
Then $\cE(G)=\cE(\phi(G,x))$ and
\begin{equation}\label{eq:e}
e(G)=n\,\cE\bigl(\phi(G,x)\bigr)-\sum_{v\in V(G)}\cE\bigl(\phi(G-v,x)\bigr).
\end{equation}
So $e(G)$ depends only on $\phi(G,x)$ and on the multiset $\{\phi(G-v,x):v\in V(G)\}$.

Let $\theta_1>\dots>\theta_d$ be the distinct eigenvalues of $A=A(G)$, with multiplicities $f_1,\dots,f_d$, and let $P_i$ be the orthogonal projection onto the eigenspace of $\theta_i$.
The numbers $\alpha_{iv}=\|P_ie_v\|$ are the \emph{angles} of $G$ \cite{CRS}.

\begin{lemma}\label{lem:angles}
For every vertex $v$ of $G$,
\[
\frac{\phi(G-v,x)}{\phi(G,x)}=\bigl((xI-A)^{-1}\bigr)_{vv}=\sum_{i=1}^d\frac{\alpha_{iv}^2}{x-\theta_i}.
\]
\end{lemma}

\begin{proof}
The polynomial $\phi(G-v,x)$ is the $(v,v)$ cofactor of $xI-A$, which gives the first equality.
The second follows from $(xI-A)^{-1}=\sum_i(x-\theta_i)^{-1}P_i$ and $(P_i)_{vv}=\|P_ie_v\|^2$.
\end{proof}

Since $\sum_v\alpha_{iv}^2=\tr P_i=f_i$, summing Lemma~\ref{lem:angles} over $v$ gives Jacobi's formula $\sum_v\phi(G-v,x)=\phi'(G,x)$.

We say that two graphs $G$ and $H$ have the \emph{same eigenvalues and angles} if $\phi(G,x)=\phi(H,x)$ and there is a bijection $\pi\colon V(G)\to V(H)$ such that $\alpha^G_{iv}=\alpha^H_{i\pi(v)}$ for all $i$ and $v$.

\begin{lemma}\label{lem:deck}
Let $\phi(G,x)=\phi(H,x)$.
Then $G$ and $H$ have the same eigenvalues and angles if and only if there is a bijection $\pi\colon V(G)\to V(H)$ with $\phi(G-v,x)=\phi(H-\pi(v),x)$ for all $v$.
In that case $e(G)=e(H)$.
\end{lemma}

\begin{proof}
By Lemma~\ref{lem:angles}, for fixed $\phi(G,x)$ the polynomial $\phi(G-v,x)$ determines the numbers $\alpha_{iv}^2$ as residues at the poles $\theta_i$, and conversely.
The last claim follows from~\eqref{eq:e}.
\end{proof}

A graph is \emph{walk-regular} if $(A^j)_{vv}$ does not depend on $v$, for every $j\ge0$ \cite{GM80}.
Vertex-transitive graphs are walk-regular.
So are strongly regular graphs: if $A^2=kI+\lambda A+\mu(J-I-A)$ and $AJ=kJ$, then every power of $A$ is a linear combination of $I$, $A$ and $J$.

\begin{lemma}\label{lem:walk}
If $G$ is walk-regular of order $n$, then $\phi(G-v,x)=\phi'(G,x)/n$ for every vertex $v$.
\end{lemma}

\begin{proof}
For $|x|$ larger than the spectral radius, $((xI-A)^{-1})_{vv}=\sum_{j\ge0}(A^j)_{vv}\,x^{-j-1}$ does not depend on $v$.
By Lemma~\ref{lem:angles}, neither does $\phi(G-v,x)$.
Jacobi's formula then gives $n\,\phi(G-v,x)=\phi'(G,x)$.
\end{proof}

\begin{corollary}\label{cor:walk}
Cospectral walk-regular graphs $G$ and $H$ of order $n$ have the same eigenvalues and angles.
Hence they are locally equienergetic, and $e(G)=e(H)=n\bigl(\cE(\phi)-\cE(\phi'/n)\bigr)$ with $\phi=\phi(G,x)$.
\end{corollary}

So every pair of non-isomorphic, cospectral, connected, walk-regular graphs is a counterexample to Conjecture~1.
Examples are pairs of non-isomorphic strongly regular graphs with the same parameters, such as $R$ and $S$ in Theorem~\ref{thm:16}, and pairs of cospectral vertex-transitive graphs (Remark~\ref{rem:twelve}).

\section{Strongly regular graphs}\label{sec:srg}

Let $G$ be a connected strongly regular graph with parameters $(n,k,\lambda,\mu)$ that is not complete.
It is well known that $A(G)$ has exactly three distinct eigenvalues $k>r>s$, with multiplicities $1,f,g$, where $r$ and $s$ are the roots of $x^2-(\lambda-\mu)x-(k-\mu)$.
Since $\tr A(G)=0$, we have $k+fr+gs=0$.

\begin{proposition}\label{prop:srg}
Let $G$ be a connected strongly regular graph of order $n$ with eigenvalues $k>r>s$ of multiplicities $1,f,g$.
Suppose that $r\ge0$ and $rs+k(fs+gr)\ge0$.
Then $\cE(G)-\cE(G-v)=2|s|$ for every vertex $v$, and $e(G)=2n|s|$.
\end{proposition}

\begin{proof}
Since $k>0$, $r\ge0$ and $k+fr+gs=0$, we have $s<0$.
By Lemma~\ref{lem:walk} and $\phi(G,x)=(x-k)(x-r)^f(x-s)^g$,
\[
\phi(G-v,x)=\frac{\phi'(G,x)}{n}=(x-r)^{f-1}(x-s)^{g-1}q(x),
\]
where $n\,q(x)=(x-r)(x-s)+f(x-k)(x-s)+g(x-k)(x-r)$.
We have $n\,q(k)>0$, $n\,q(r)=f(r-k)(r-s)<0$ and $n\,q(s)=g(s-k)(s-r)>0$.
So $q$ has a root $\rho_1\in(r,k)$ and a root $\rho_2\in(s,r)$, and $\rho_1>0$.
Moreover $n\,q(0)=rs+k(fs+gr)\ge0$.
If $\rho_2<0$, then $0$ would lie strictly between the roots of $q$ and $q(0)<0$.
Hence $\rho_2\ge0$.
The trace of $A(G-v)$ is zero, so $(f-1)r+(g-1)s+\rho_1+\rho_2=0$.
Together with $k+fr+gs=0$ this gives $\rho_1+\rho_2=k+r+s$.
Therefore
\[
\cE(G-v)=(f-1)r+(g-1)|s|+\rho_1+\rho_2=(f-1)r+(g-1)|s|+k+r+s,
\]
while $\cE(G)=k+fr+g|s|$.
Subtracting gives $\cE(G)-\cE(G-v)=|s|-s=2|s|$.
\end{proof}

\begin{remark}
The sign condition cannot be omitted.
For the Petersen graph ($k=3$, $r=1$, $s=-2$, $f=5$, $g=4$) we have $rs+k(fs+gr)=-20$.
Here $\phi(G-v,x)=(x-1)^4(x+2)^3(x^2-2x-2)$, and the roots $1\pm\sqrt3$ of the last factor have opposite signs.
This gives $e(G)=10\bigl(16-(10+2\sqrt3)\bigr)=60-20\sqrt3$.
\end{remark}

The lattice graph $L_2(\ell)=K_\ell\square K_\ell$ has eigenvalues $2\ell-2$, $\ell-2$, $-2$ with multiplicities $1$, $2(\ell-1)$, $(\ell-1)^2$.
For it,
\[
rs+k(fs+gr)=-2(\ell-2)+2(\ell-1)^2\bigl((\ell-1)(\ell-2)-4\bigr),
\]
which equals $32$ for $\ell=4$ and is positive for all $\ell\ge4$.
The triangular graph $T(m)=L(K_m)$ has eigenvalues $2m-4$, $m-4$, $-2$ with multiplicities $1$, $m-1$, $m(m-3)/2$.
For it,
\[
rs+k(fs+gr)=-2(m-4)+2(m-2)\Bigl(\frac{m(m-3)(m-4)}2-2(m-1)\Bigr),
\]
which equals $60$ for $m=6$ and is positive for all $m\ge6$.
Proposition~\ref{prop:srg} gives
\[
e\bigl(L_2(\ell)\bigr)=4\ell^2\quad(\ell\ge4),\qquad e\bigl(T(m)\bigr)=2m(m-1)\quad(m\ge6).
\]

\begin{corollary}\label{cor:fam}
Let $m\ge6$ and $\ell\ge4$ satisfy $m(m-1)/2=\ell^2$.
Then $T(m)$ and $L_2(\ell)$ are connected graphs of order $\ell^2$, and $e(T(m))=e(L_2(\ell))=4\ell^2$.
Their degrees $2m-4$ and $2\ell-2$ are different, so they are not cospectral.
The equation $m(m-1)/2=\ell^2$ is equivalent to $(2m-1)^2-8\ell^2=1$ and has infinitely many solutions, beginning with $(m,\ell)=(9,6),(50,35),(289,204)$.
\end{corollary}

\begin{proof}
Only the claim about the degrees needs comment.
If $2m-4=2\ell-2$, then $m=\ell+1$ and $(\ell+1)\ell/2=\ell^2$, so $\ell=1$.
\end{proof}

For example, $T(9)$ and $L_2(6)$ have $36$ vertices, degrees $14$ and $10$, and local energy $144$.
We confirmed $e(T(9))=e(L_2(6))=144$ and $e(T(50))=e(L_2(35))=4900$ numerically.

\section{Three graphs of order 16}\label{sec:sixteen}

Let $\mathbb{Z}_4^2=\mathbb{Z}_4\times\mathbb{Z}_4$.
For a set $C\subseteq\mathbb{Z}_4^2\setminus\{0\}$ with $C=-C$, the Cayley graph $\Cay(\mathbb{Z}_4^2,C)$ has vertex set $\mathbb{Z}_4^2$, and $x\sim y$ if and only if $x-y\in C$.
Cayley graphs are vertex-transitive.
We use
\begin{align*}
C_R&=\{(0,\pm1),(0,2),(\pm1,0),(2,0)\},\\
C_S&=\{\pm(1,0),\pm(0,1),\pm(1,1)\},\\
C_Q&=C_R\cup\{(\pm1,\pm1)\},
\end{align*}
and put $R=\Cay(\mathbb{Z}_4^2,C_R)$, $S=\Cay(\mathbb{Z}_4^2,C_S)$ and $Q=\Cay(\mathbb{Z}_4^2,C_Q)$.
In $R$, two distinct vertices are adjacent if and only if they agree in one coordinate, so $R=K_4\square K_4=L_2(4)$.
The graph $S$ is the Shrikhande graph \cite{Shr}.
In a Cayley graph the number of common neighbours of $x$ and $y$ is the number of $c\in C$ with $x-y-c\in C$.
A direct count shows that $S$ is strongly regular with parameters $(16,6,2,2)$ and that $Q$ is strongly regular with parameters $(16,10,6,6)$.
We checked by computer that $Q$ is isomorphic to the complement of the Clebsch graph.
Here the Clebsch graph is the Cayley graph of $\mathbb{Z}_2^4$ whose connection set consists of the four unit vectors and the all-ones vector.

\begin{proof}[Proof of Theorem~\ref{thm:16}]
The data are collected in Table~\ref{tab:16}.
The graphs $R$ and $S$ have eigenvalues $6,2,-2$ with multiplicities $1,6,9$.
So $\cE(R)=\cE(S)=6+12+18=36$, and $rs+k(fs+gr)=-4+6(-12+18)=32$.
The graph $Q$ has eigenvalues $10,2,-2$ with multiplicities $1,5,10$.
So $\cE(Q)=10+10+20=40$, and $rs+k(fs+gr)=-4+10(-10+20)=96$.
By Proposition~\ref{prop:srg}, $e(R)=e(S)=e(Q)=2\cdot16\cdot2=64$.
The deleted graphs have
\[
\phi(R-v)=\phi(S-v)=(x-2)^5(x+2)^8(x^2-6x+2),\qquad \phi(Q-v)=(x-2)^4(x+2)^9(x^2-10x+6),
\]
whose quadratic factors have the positive roots $3\pm\sqrt7$ and $5\pm\sqrt{19}$.
Strongly regular graphs with $\mu>0$ have diameter $2$, so the three graphs are connected.
Each row of $R$ induces a $K_4$.
In $S$, the neighbourhood $C_S$ of $0$ induces the $6$-cycle
\[
(1,0),\ (1,1),\ (0,1),\ (3,0),\ (3,3),\ (0,3),
\]
which has no triangle.
So no $K_4$ of $S$ contains $0$, and since $S$ is vertex-transitive, $S$ contains no $K_4$.
Hence $R\not\cong S$.
Finally, $Q$ is $10$-regular, while $R$ and $S$ are $6$-regular.
\end{proof}

\begin{table}[ht]
\caption{The graphs of Theorem~\ref{thm:16}. Here $v$ is any vertex.}\label{tab:16}
\begin{tabular}{lccccc}
\toprule
graph & parameters & spectrum & $\cE(G)$ & $\cE(G-v)$ & $e(G)$\\
\midrule
$R=K_4\square K_4$ & $(16,6,2,2)$ & $6^1,\ 2^6,\ (-2)^9$ & $36$ & $32$ & $64$\\
$S$ (Shrikhande) & $(16,6,2,2)$ & $6^1,\ 2^6,\ (-2)^9$ & $36$ & $32$ & $64$\\
$Q$ (co-Clebsch) & $(16,10,6,6)$ & $10^1,\ 2^5,\ (-2)^{10}$ & $40$ & $36$ & $64$\\
\bottomrule
\end{tabular}
\end{table}

The equality $e(R)=e(S)$ also has a direct certificate.
Let $T=\{(0,1),(1,0),(2,2),(3,3)\}$.
It is an independent set of $R$, and every vertex outside $T$ has exactly two neighbours in $T$, namely the elements of $T$ in its row and in its column.
Godsil--McKay switching \cite{GM82} with respect to $T$ therefore turns $R$ into a cospectral graph, and this graph is isomorphic to $S$.
The switching matrix is $M=\mathrm{diag}(\tfrac12J-I,\,I)$, which is orthogonal.
Composing $2M$ with the permutation matrix of an isomorphism onto $S$ that fixes $(0,0)$ gives an integer matrix $P$ with $PA(R)=A(S)P$ and $PP^{\mathsf T}=4I$.
Since $(0,0)\notin T$, the matrix $P$ maps $e_{(0,0)}$ to $2e_{(0,0)}$.
So deleting the row and the column of $(0,0)$ gives a certificate of the same kind for $R-(0,0)$ and $S-(0,0)$.
This is the certificate used in the formal proof.

\section{The exhaustive search up to order 10}\label{sec:search}

We computed the local energy of every connected graph of order at most $10$.
The graphs were generated by \texttt{geng -c} from nauty \cite{MP}, one graph per isomorphism class.
Table~\ref{tab:search} lists the numbers of graphs.
The method has three steps.
\begin{enumerate}
\item \emph{Floating point.} For each graph, $e(G)$ was computed in double precision with a symmetric eigenvalue solver, using~\eqref{eq:e}.
The values were sorted, and consecutive values closer than a tolerance $\tau$ were put into the same group of candidates; we used $\tau=10^{-8}$ for $n\le9$ and $\tau=10^{-10}$ for $n=10$.
\item \emph{Exact equality.} For each pair of graphs in a group, the integer characteristic polynomials $\phi(G,x)$ and the multisets $\{\phi(G-v,x)\}$ were computed exactly by the Faddeev--LeVerrier algorithm.
If they agree for both graphs, the local energies are equal by Lemma~\ref{lem:deck}.
\item \emph{Inequality.} Otherwise, both local energies were recomputed with $30$ significant digits.
A pair is declared unequal if the difference exceeds $10^{-20}$.
\end{enumerate}
The only pairs that were equal but not settled by step~2 are the pairs $(K_n,C_n)$, whose local energies are equal by the formulas of Espinal and Rada.

For $n\le8$ the only candidate groups are the pairs $(K_5,C_5)$ and $(K_7,C_7)$.
For $n=9$ there are $112$ groups, and steps~2 and~3 settle all of them: the pair $(K_9,C_9)$, whose $30$-digit values differ by about $10^{-29}$, and $111$ unequal pairs, the closest of which differ by about $1.9\cdot10^{-10}$.

For $n=10$ there are $2187$ groups.
Steps~2 and~3 give exactly $58$ pairs with the same eigenvalues and angles, and $2131$ unequal pairs.
The smallest difference between the local energies of an unequal pair is about $1.3\cdot10^{-14}$.

Only step~1 is not exact.
For the $58$ pairs with equal local energy, the double-precision values differ by at most $1.5\cdot10^{-13}$, three orders of magnitude below $\tau$.
For pairs in different groups, the computed values differ by at least $\tau$.

We checked the following properties of the $58$ pairs independently of the search.
In each pair the graphs are connected and non-isomorphic, and they have the same characteristic polynomial and the same multiset of vertex-deleted characteristic polynomials.
They also have the same degree sequence and cospectral complements.
In each pair, one graph is isomorphic to the graph obtained from the other by Godsil--McKay switching with respect to four vertices.
The numbers of edges range from $16$ to $29$, and exactly one pair has $16$ edges.
This proves Theorem~\ref{thm:10}, up to the floating-point step.
The programs are \texttt{order10/le\_scan.py}, \texttt{order10/exact\_pairs.py}, \texttt{order10/exact\_check.py} and \texttt{order10/verify58.py}.
The list of the $58$ pairs, in graph6 format, is \texttt{order10/pairs10.txt}.

\begin{table}[ht]
\caption{Connected graphs of order $n$ and their pairs of non-isomorphic locally equienergetic graphs.}\label{tab:search}
\begin{tabular}{rrl}
\toprule
$n$ & connected graphs & locally equienergetic pairs\\
\midrule
$\le4$ & $10$ & none\\
$5$ & $21$ & $(K_5,C_5)$\\
$6$ & $112$ & none\\
$7$ & $853$ & $(K_7,C_7)$\\
$8$ & $11117$ & none\\
$9$ & $261080$ & $(K_9,C_9)$\\
$10$ & $11716571$ & $58$ pairs, all with the same eigenvalues and angles\\
\bottomrule
\end{tabular}
\end{table}

\section{The pair of order 10}\label{sec:ten}

Let $G$ and $H$ be the graphs on the vertex set $\{0,1,\dots,9\}$ with the edges listed in Table~\ref{tab:ten}.
In graph6 format they are \verb|I?qa``kWo| and \verb|ICOeeOwFG|.
Both have $16$ edges and degree sequence $(4,4,4,4,3,3,3,3,2,2)$, and both are connected.
Their characteristic polynomial is
\[
\phi(G,x)=\phi(H,x)=x^2(x-2)(x+2)^3(x^2-2)(x^2-4x+2),
\]
and $\cE(G)=\cE(H)=12+2\sqrt2$.
Their local energy is $e(G)=e(H)=21.2203488379522\ldots$

\begin{table}[ht]
\caption{The pair of order $10$ with the fewest edges.}\label{tab:ten}
\begin{tabular}{ll}
\toprule
graph & edges\\
\midrule
$G$ & 04 05 14 16 18 19 25 27 28 29 36 37 48 58 69 79\\
$H$ & 03 06 07 14 16 17 25 28 36 38 39 47 48 49 59 89\\
\bottomrule
\end{tabular}
\end{table}

\begin{proposition}\label{prop:ten}
The graphs $G$ and $H$ are connected, non-isomorphic and locally equienergetic.
More precisely, let $\pi$ be the permutation $(0\,2\,4)(1\,3\,5)(6\,7)(8\,9)$ of $\{0,\dots,9\}$.
Then $\phi(G-v,x)=\phi(H-\pi(v),x)$ for every vertex $v$.
\end{proposition}

\begin{proof}
The sets $C_1=\{0,6,7,8\}$ and $C_2=\{3,4,5,9\}$ are independent in $G$.
Every vertex outside $C_i$ has exactly two neighbours in $C_i$.
Let $G_i$ be the graph obtained from $G$ by Godsil--McKay switching with respect to $C_i$.
This means that every vertex outside $C_i$ exchanges its two neighbours in $C_i$ for the other two vertices of $C_i$.
By \cite{GM82}, $G_i$ is cospectral with $G$.
A computer check shows that $G_1$ and $G_2$ are both isomorphic to $H$.

Let $v\notin C_i$.
Deleting $v$ preserves the switching conditions, and switching $G-v$ with respect to $C_i$ gives $G_i-v$.
So $G-v$ is cospectral with $G_i-v$, which is isomorphic to $H-\sigma_i(v)$ for an isomorphism $\sigma_i\colon G_i\to H$.
Every vertex lies outside $C_1$ or outside $C_2$, because $C_1\cap C_2=\emptyset$.
The isomorphisms $\sigma_1$ and $\sigma_2$ can be chosen so that $\sigma_2$ on $C_1$ and $\sigma_1$ on the other six vertices together give the permutation $\pi$.
This gives $\phi(G-v,x)=\phi(H-\pi(v),x)$ for all $v$, and $e(G)=e(H)$ by Lemma~\ref{lem:deck}.

The graph $H$ contains the triangles $\{8,9,3\}$ and $\{8,9,4\}$, which share the edge $\{8,9\}$.
The four triangles of $G$ are $\{1,4,8\}$, $\{1,6,9\}$, $\{2,7,9\}$ and $\{2,5,8\}$, and any two of them share at most one vertex.
Hence $G\not\cong H$.
The diameters are also different: $G$ has diameter $4$ and $H$ has diameter $3$.
\end{proof}

Since $10$ is even, the pair $(G,H)$ is not of the form $(K_n,C_n)$ with $n$ odd.
So it is a counterexample both to Corollary~2 and to Conjecture~1 of \cite{DP}.
The switching argument also provides the certificates for the formal proof.
Let $M_i=\mathrm{diag}(\tfrac12J-I,\,I)$ be the switching matrix for $C_i$.
Then $P_i=2\Pi_{\sigma_i}M_i$ is an integer matrix with entries in $\{0,\pm1,2\}$, where $\Pi_{\sigma_i}$ is the permutation matrix of $\sigma_i$.
It satisfies $P_iA(G)=A(H)P_i$ and $P_iP_i^{\mathsf T}=4I$.
For $v\notin C_i$, deleting row $\pi(v)$ and column $v$ of $P_i$ gives a matrix $P$ with $PA(G-v)=A(H-\pi(v))P$ and $PP^{\mathsf T}=4I$.

\begin{remark}\label{rem:twelve}
There are also cospectral vertex-transitive counterexamples of order $12$.
Let $C_{12}(1,2)$ be the Cayley graph of $\mathbb{Z}_{12}$ with connection set $\{\pm1,\pm2\}$.
Let $D_6=\langle r,s\mid r^6=s^2=1,\ srs=r^{-1}\rangle$ be the dihedral group of order $12$.
Let $X=\Cay(D_6,\{s,rs,r^3,r^3s\})$, in which $g$ is adjacent to $gt$ for each $t$ in the connection set.
The graph $X$ consists of three copies of $K_4$, on the cosets of the subgroup $\{1,r^3,s,r^3s\}$, joined by the perfect matching $g\mapsto g\,rs$.
Both graphs are connected and $4$-regular, and both have characteristic polynomial
\[
x^3(x-4)(x+2)^4(x^2-2x-2)^2.
\]
They are not isomorphic, because $X$ contains $K_4$ and $C_{12}(1,2)$ does not.
By Corollary~\ref{cor:walk} they are locally equienergetic, with $e=23.6933935\ldots$
We found these graphs in a numerical survey of the connected Cayley graphs of all groups of order at most $15$.
In that survey, the only other coincidences of local energy were the pairs $(K_n,C_n)$ for odd $n$ and three further cospectral pairs of order $12$.
\end{remark}

\section{Formal verification}\label{sec:lean}

The definitions, the statements and the counterexamples were formalized in Lean~4 (version 4.33.1) with Mathlib.
The energy is defined through the eigenvalues of the real adjacency matrix, which is Hermitian.
The graph $G-v$ is the subgraph induced on the complement of $\{v\}$.
Figure~\ref{fig:defs} shows the definitions.

\begin{figure}[t]
\begin{Verbatim}[fontsize=\small,frame=single]
noncomputable def energy {V : Type*} [Fintype V] (G : SimpleGraph V) : ℝ := by
  classical
  exact ∑ i, |(G.isHermitian_adjMatrix ℝ).eigenvalues i|

noncomputable def localVertexEnergy {V : Type*} [Fintype V]
    (G : SimpleGraph V) (v : V) : ℝ := by
  classical
  exact energy G - energy (G.induce ({v}ᶜ : Set V))

noncomputable def localEnergy {V : Type*} [Fintype V]
    (G : SimpleGraph V) : ℝ :=
  ∑ v, localVertexEnergy G v
\end{Verbatim}
\caption{The definitions of energy and local energy.}\label{fig:defs}
\end{figure}

Two graphs are locally equienergetic if they have the same number of vertices and the same local energy.
The proposition \texttt{Conjecture1} says the following.
For all finite vertex types $V_1,V_2$ and graphs $G_1,G_2$, if $G_1$ and $G_2$ are connected, not isomorphic and locally equienergetic, then the order $n$ is odd and the pair is isomorphic to $(K_n,C_n)$ or to $(C_n,K_n)$.
Here $K_n$ is the complete graph on \texttt{Fin n} and $C_n$ is Mathlib's \texttt{cycleGraph n}.
The proposition \texttt{Corollary2} adds the hypothesis $n\le10$.
The proposition \texttt{Conjecture1NonEquienergetic} replaces the hypothesis that the graphs are not isomorphic by the hypothesis that their energies differ.
The statement file \texttt{ComparatorChallenges/LocalEnergy.lean} of the repository contains these definitions and the three statements \mbox{\texttt{¬ Conjecture1}}, \mbox{\texttt{¬ Conjecture1NonEquienergetic}} and \mbox{\texttt{¬ Corollary2}}.
The proof file \texttt{KOS/ChemicalGraphTheory/LocalEnergy/Main.lean} begins with the same definitions, copied verbatim, and proves all three statements; Comparator confirms that it proves exactly the statements of the statement file.

The theorem \mbox{\texttt{not\_conjecture1}} is proved with the pair $(R,S)$.
The proof file also proves the following:
\begin{itemize}
\item the values $\cE=36,36,40$, $\cE(G-v)=32,32,36$ and $e=64$ for $R$, $S$ and $Q$;
\item the theorem \texttt{not\_conjecture1NonEquienergetic}, which refutes the variant of Conjecture~1 that only considers pairs of graphs with different energies, with the pair $(R,Q)$.
\end{itemize}

For the pair of order $10$ it proves three theorems:
\begin{itemize}
\item \texttt{order10\_counterexample}: $G$ and $H$ are connected, not isomorphic and have equal local energy;
\item \texttt{not\_corollary2};
\item \texttt{not\_conjecture1\_order10}.
\end{itemize}

The proofs reduce everything to finite computations in four steps.
\begin{itemize}
\item Mathlib's description of the characteristic polynomial of a Hermitian matrix shows that the energy is the sum of the absolute values of the roots of the characteristic polynomial.
\item An integer matrix $P$ with $PA=BP$ and $PP^{\mathsf T}=4I$ then shows that two graphs have the same energy.
\item Vertex-deleted subgraphs are identified with graphs on \texttt{Fin 15} or \texttt{Fin 9}.
\item All finite facts are checked by the kernel with \texttt{decide}: the certificate identities, connectivity, the forbidden subgraphs used for non-isomorphism, and the integer identities and traces used for the exact energies.
\end{itemize}
The exact values use a lemma of the following form: if $p(A)=0$ and $g$ agrees with $|x|$ on the roots of $p$, then $\cE=\tr g(A)$.
The proof of that lemma uses the continuous functional calculus in Mathlib.
For example, $(A-6)(A-2)(A+2)=0$ for $A=A(R)$ and $g(x)=\tfrac32+\tfrac{x^2}8$.

Checking the proof file takes about one minute on one core.
All main theorems depend only on the axioms \texttt{propext}, \texttt{Classical.choice} and \texttt{Quot.sound}.
The general results of Sections~\ref{sec:angles} and~\ref{sec:srg}, the search of Section~\ref{sec:search} and Remark~\ref{rem:twelve} are not formalized.

\subsection*{Author contributions}
GPT-6-Astra, acting as an advisor surveying open problems, identified the conjecture and proposed the pair $(R,S)$, which it checked in exact arithmetic.
Claude Opus~5.5 did the following:
\begin{itemize}
\item verified the pair $(R,S)$;
\item found the pair $(R,Q)$ and the general mechanisms of Sections~\ref{sec:angles} and~\ref{sec:srg};
\item formalized the results in Lean;
\item carried out the exhaustive search for $n\le10$;
\item formalized the pair of order $10$;
\item wrote the manuscript.
\end{itemize}
Claude Fable~5.1 refereed the proofs and the formalization, and verified the spectra and the example of order $12$ independently.
A human operator initiated the project and provided computing resources; the mathematics, the formal proofs and the text were produced by the models.

\begin{thebibliography}{9}

\bibitem{CRS}
D.~Cvetkovi\'c, P.~Rowlinson and S.~Simi\'c,
\emph{Eigenspaces of Graphs},
Cambridge University Press, Cambridge, 1997.

\bibitem{DP}
C.~Dede and K.~M.~Popat,
\emph{Locally equienergetic graphs},
MATCH Commun. Math. Comput. Chem. \textbf{93} (2025), no.~3, 759--766; doi:10.46793/match.93-3.759D; arXiv:2604.14686.

\bibitem{ER}
C.~Espinal and J.~Rada,
\emph{Graph energy change due to vertex deletion},
MATCH Commun. Math. Comput. Chem. \textbf{92} (2024), no.~1, 89--103; doi:10.46793/match.92-1.089E.

\bibitem{GM80}
C.~D.~Godsil and B.~D.~McKay,
\emph{Feasibility conditions for the existence of walk-regular graphs},
Linear Algebra Appl. \textbf{30} (1980), 51--61.

\bibitem{GM82}
C.~D.~Godsil and B.~D.~McKay,
\emph{Constructing cospectral graphs},
Aequationes Math. \textbf{25} (1982), 257--268.

\bibitem{LSG}
X.~Li, Y.~Shi and I.~Gutman,
\emph{Graph Energy},
Springer, New York, 2012.

\bibitem{MP}
B.~D.~McKay and A.~Piperno,
\emph{Practical graph isomorphism, II},
J. Symbolic Comput. \textbf{60} (2014), 94--112.

\bibitem{Shr}
S.~S.~Shrikhande,
\emph{The uniqueness of the $L_2$ association scheme},
Ann. Math. Statist. \textbf{30} (1959), 781--798.

\end{thebibliography}

\end{document}
